\documentclass[border=12pt]{standalone}
\usepackage[no-math]{fontspec}
\usepackage[brazil]{babel}
\usepackage{amsmath,amssymb}
\IfFileExists{C:/Windows/Fonts/Lato_Regular.ttf}{%
  \setmainfont{Lato}[
    Path = C:/Windows/Fonts/,
    UprightFont = Lato_Regular.ttf,
    BoldFont = Lato_Bold.ttf,
    ItalicFont = Lato_Italic.ttf,
    BoldItalicFont = Lato_BoldItalic.ttf
  ]
  \setsansfont{Lato}[
    Path = C:/Windows/Fonts/,
    UprightFont = Lato_Regular.ttf,
    BoldFont = Lato_Bold.ttf,
    ItalicFont = Lato_Italic.ttf,
    BoldItalicFont = Lato_BoldItalic.ttf
  ]
}{%
  \RequirePackage[default]{lato}
}
\renewcommand{\familydefault}{\sfdefault}

\usepackage{xcolor}
\usepackage{tikz}
\usepackage{pgfplots}
\pgfplotsset{compat=1.18}
\usetikzlibrary{arrows.meta,positioning,shapes.geometric,fit,backgrounds,calc,decorations.pathreplacing,patterns}

% Paleta institucional UEFS / Extensão Rural
\definecolor{erNavy}{HTML}{1F4268}
\definecolor{erNavyDark}{HTML}{152E4A}
\definecolor{erGold}{HTML}{B38A3D}
\definecolor{erGoldDark}{HTML}{8C6927}
\definecolor{erTeal}{HTML}{0E7C7B}
\definecolor{erForest}{HTML}{2D6A4F}
\definecolor{erAlert}{HTML}{9E2A2B}
\definecolor{erInk}{HTML}{26313D}
\definecolor{erMuted}{HTML}{5B6573}

\definecolor{bgNavy}{HTML}{EDF2F7}
\definecolor{bgGold}{HTML}{F8F3E9}
\definecolor{bgTeal}{HTML}{EBF6F5}
\definecolor{bgForest}{HTML}{EEF5F0}
\definecolor{bgAlert}{HTML}{FDF2F2}

\tikzset{
  >=Stealth,
  erBox/.style={
    draw=#1, fill=white, rounded corners=5pt, line width=1.15pt,
    inner sep=7pt, align=center, font=\footnotesize\sffamily, text=erInk
  },
  erTint/.style={
    draw=#1, fill=#1!8!white, rounded corners=5pt, line width=1.15pt,
    inner sep=7pt, align=center, font=\footnotesize\sffamily, text=erInk
  },
  erBadge/.style={
    fill=#1, text=white, rounded corners=9pt, inner xsep=8pt, inner ysep=3pt,
    font=\scriptsize\bfseries\sffamily
  },
  erTitle/.style={font=\normalsize\bfseries\sffamily, text=erNavy, align=center},
  erSub/.style={font=\footnotesize\sffamily, text=erMuted, align=center},
  erArrow/.style={->, >=Stealth, line width=1.35pt},
  erDArrow/.style={<->, >=Stealth, line width=1.35pt}
}


\begin{document}
\begin{tikzpicture}

  % Título principal
  \node[erTitle, font=\large\bfseries\sffamily] (title) at (0, 5.2)
    {Arquitetura Institucional da ATER no Brasil: Pública, Terceiro Setor e Privada};

  % ================= 3 CARDS BALANCEADOS =================
  % Card Esquerdo: ATER Pública & Governamental
  \node[draw=erForest, fill=bgForest, rounded corners=7pt, line width=1.2pt,
        minimum width=5.15cm, minimum height=5.9cm] (card1) at (-5.55, 1.45) {};

  \node[erBadge=erForest, rounded corners=5pt, inner ysep=4pt, minimum width=4.55cm, align=center]
    at ([yshift=-20pt]card1.north) {ATER Pública \& Governamental};

  \node[erBox=erForest!60, minimum width=4.65cm, minimum height=4.35cm, text width=4.15cm, align=left]
    at ([yshift=-16pt]card1.center) {
    {\color{erForest}\textbf{\scriptsize Governança \& Atores}}\\[1pt]
    \scriptsize EMATERs, BAHIATER, INCRA e Secretarias Municipais.\\[7pt]
    {\color{erForest}\textbf{\scriptsize Financiamento}}\\[1pt]
    \scriptsize Orçamento público federal, estadual e municipal (gratuidade).\\[7pt]
    {\color{erForest}\textbf{\scriptsize Orientação Estratégica}}\\[1pt]
    \scriptsize Foco em Agricultura Familiar e Política Pública (\textbf{Lei 12.188/10}).
  };

  % Card Central: ONGs, Cooperativas & Terceiro Setor
  \node[draw=erTeal, fill=bgTeal, rounded corners=7pt, line width=1.2pt,
        minimum width=5.15cm, minimum height=5.9cm] (card2) at (0, 1.45) {};

  \node[erBadge=erTeal, rounded corners=5pt, inner ysep=4pt, minimum width=4.55cm, align=center]
    at ([yshift=-20pt]card2.north) {ONGs, Cooperativas \& 3º Setor};

  \node[erBox=erTeal!60, minimum width=4.65cm, minimum height=4.35cm, text width=4.15cm, align=left]
    at ([yshift=-16pt]card2.center) {
    {\color{erTeal}\textbf{\scriptsize Governança \& Atores}}\\[1pt]
    \scriptsize ASA, COOPAFS, sindicatos e Redes Agroecológicas.\\[7pt]
    {\color{erTeal}\textbf{\scriptsize Financiamento}}\\[1pt]
    \scriptsize Chamadas Públicas da PNATER, convênios e fundos territoriais.\\[7pt]
    {\color{erTeal}\textbf{\scriptsize Orientação Estratégica}}\\[1pt]
    \scriptsize Enraizamento territorial, convivência com o semiárido e autonomia.
  };

  % Card Direito: ATER Privada, Integradoras & Sistema S
  \node[draw=erGold, fill=bgGold, rounded corners=7pt, line width=1.2pt,
        minimum width=5.15cm, minimum height=5.9cm] (card3) at (5.55, 1.45) {};

  \node[erBadge=erGoldDark, rounded corners=5pt, inner ysep=4pt, minimum width=4.55cm, align=center]
    at ([yshift=-20pt]card3.north) {ATER Privada, Integradoras \& Sist. S};

  \node[erBox=erGold!70, minimum width=4.65cm, minimum height=4.35cm, text width=4.15cm, align=left]
    at ([yshift=-16pt]card3.center) {
    {\color{erGoldDark}\textbf{\scriptsize Governança \& Atores}}\\[1pt]
    \scriptsize Revendas de insumos, Agroindústrias e SENAR / ATeG.\\[7pt]
    {\color{erGoldDark}\textbf{\scriptsize Financiamento}}\\[1pt]
    \scriptsize Margem comercial, contratos de integração ou contribuição setorial.\\[7pt]
    {\color{erGoldDark}\textbf{\scriptsize Orientação Estratégica}}\\[1pt]
    \scriptsize Foco comercial, pacotes tecnológicos ou eficiência na cadeia produtiva.
  };

  % ================= BANNER SÍNTESE INFERIOR =================
  \node[draw=erNavy, fill=bgNavy, rounded corners=6pt, line width=1.2pt,
        minimum width=16.25cm, inner xsep=12pt, inner ysep=9pt, align=center,
        font=\footnotesize\sffamily, text=erNavy] (banner) at (0, -2.35) {
    \textbf{Desafio ao Engenheiro Agrônomo:}
    \textit{prevenir conflitos de interesse técnico-comerciais e assegurar a autonomia decisória do agricultor.}
  };

\end{tikzpicture}
\end{document}
